\documentclass[11pt]{article}

\usepackage[margin=1in]{geometry}
\usepackage[utf8]{inputenc}
\usepackage[T1]{fontenc}
\usepackage{booktabs}
\usepackage{natbib}
\usepackage{hyperref}
\usepackage{graphicx}
\usepackage{amsmath}
\usepackage{verbatim}
\usepackage{booktabs}
\usepackage{array}
\usepackage{xcolor}
\usepackage{microtype}
\usepackage{multicol}
\usepackage{enumitem}
\usepackage{ulem}
\usepackage{tikz}
\usetikzlibrary{arrows.meta,positioning,calc,backgrounds,shapes.geometric,fit}
%\usepackage[font={small,sf},labelfont={bf,sf}]{caption}
\usepackage{titlesec}
\usepackage{mathptmx}
\usepackage[misc]{ifsym}
\usepackage[scaled=0.92]{helvet}
\usepackage[margin=1in]{geometry}

\definecolor{sciblue}{RGB}{0,84,147}
\definecolor{scipale}{RGB}{235,242,248}
\definecolor{apipastel}{RGB}{240,168,175}   % pastel rose fill
\definecolor{apidark}{RGB}{168,66,78}       % readable rose for text
\definecolor{apired}{RGB}{176,42,55}
\definecolor{locpastel}{RGB}{154,208,192}   % pastel green fill
\definecolor{locdark}{RGB}{29,113,97}       % readable green for text
\definecolor{localteal}{RGB}{0,120,110}
\definecolor{goldpastel}{RGB}{246,224,160}  % pastel amber fill
\definecolor{golddark}{RGB}{150,110,20}     % readable amber for text
\definecolor{roundgold}{RGB}{176,124,10}

\newcommand{\florian}[1]{\textcolor{apidark}{florian: #1}}


\newcolumntype{L}[1]{>{\raggedright\arraybackslash}p{#1}}
\newcolumntype{C}[1]{>{\centering\arraybackslash}p{#1}}

\bibpunct{(}{)}{;}{a}{}{,}
\usepackage{xcolor}
\hypersetup{
citebordercolor=white,
filebordercolor=red,
linkbordercolor=blue
}
\hypersetup{
colorlinks,
citecolor=blue,
linkcolor=red,
urlcolor=blue}

\usepackage{palatino}
\fontfamily{ppl}


\title{Randomness in large language models: \vspace{1mm} \\ What researchers %in economics and management\\
need to know (and report)%\thanks{We would like to thank [xxx] for their helpful comments and suggestions. The replication package of this paper is available at [URL here].}
\\\vspace{4mm}}

\author{
Guillaume Coqueret\thanks{EMLYON Business School, France. \Letter  \hspace{1mm} coqueret@em-lyon.com. Coqueret is affiliated with Cascad, a nonprofit certification agency that verifies the reproducibility of research articles in economics and management.} \hspace{6mm}  \and
Joan Llull\thanks{Institute for Economic Analysis (IAE-CSIC), and Barcelona School of Economics, Spain. \Letter \hspace{1mm} joan.llull@iae.csic.es. Llull is the Data Editor of the Econometric Society.}  \hspace{6mm}  \and
Florian Oswald\thanks{University of Turin and Collegio Carlo Alberto, Italy. \Letter \hspace{1mm} florian.oswald@unito.it. Oswald is the Data Editor of the Journal of Political Economy.} \and  
Christophe P\'erignon\thanks{HEC Paris, France. \Letter  \hspace{1mm} perignong@hec.fr. P\'erignon is co-founder and Director of Cascad.}   \and  
Christoph Scheuch\thanks{Humboldt University of Berlin, Germany. \Letter \hspace{1mm} christoph.scheuch@hu-berlin.de. Scheuch is the Data Editor of the Review of Financial Studies and is affiliated with Cascad.}     \and
Lars Vilhuber\thanks{Cornell University, United States. \Letter \hspace{1mm} lars.vilhuber@cornell.edu. Vilhuber is the Data Editor of the American Economic Association.}
}

\date{}
\begin{document}
\maketitle

\begin{abstract}
% i deleted "for empirical research in economics and management science" . 
Large language models (LLMs) are increasingly used to generate data for research. Typical use cases are classifications, annotations, information extraction, and generation of numerical scores. Unlike conventional measurements, LLM outputs can vary across repeated requests even when the prompt and apparent model settings remain unchanged. This variation arises from deliberate sampling, silent model updates, numerical rounding, or expert routing. Setting a dedicated \textit{temperature} parameter to zero removes deliberate sampling when that option is available, but it does not eliminate the other sources of randomness. Exact reproduction is therefore generally not possible when using proprietary application programming interfaces. Local execution of open-weight models offers greater control, but reproducibility still depends on the complete hardware and software stack. We illustrate these issues through sentiment classifications of corporate filings and examine their consequences for downstream regression results. We then propose a reporting standard for articles and replication packages, as well as guidance for data editors and authors. Together, these findings and recommendations establish that LLM outputs should be treated as draws from a distribution rather than as fixed measurements.
\end{abstract}

\section{Introduction}

Large language models (LLMs) have become ubiquitous in modern academic research. Statements such as ``\textit{we use a large language model to classify 60,000 documents}'' now appear routinely in empirical studies in economics and management. %In a \href{https://nber-poll.paulgp.workers.dev/results}{July 2026 NBER survey of economists}, 60\% of them reported that LLMs are an essential part of their research workflow. 
Researchers use these models to classify and annotate corporate filings \citep{wang2025assessing} and earnings announcements \citep{siano2025news,harford2026}, measure sentiment or identify topics in social media posts and news articles \citep{gilardi2023}, interpret and summarize central bank communications \citep{hansen2023}, extract structured information from historical records and images \citep{dell2025}, predict financial outcomes \citep{lopezlira2025}, and generate customer recommendations \citep{kweon2025uncertainty}. The resulting model outputs increasingly serve as inputs for subsequent analyses.%deleted empirical 
\medskip

Unlike traditional measurements, however, LLM-generated data are generally \textit{not deterministic} \citep{ouyang2024,he2025,deKokMS2025}. The same prompt may yield different outputs, and therefore different data, even when the prompt and apparent model settings are unchanged. For commercial providers such as OpenAI or Anthropic, researchers typically cannot observe the serving infrastructure, archive the deployed model, identify the hardware on which it ran, or even control the \textit{temperature} -- a parameter akin to a logit scale, determining the amount of randomness in model output -- in every interface. These hidden sources of variation are rarely acknowledged, and the resulting data are typically analyzed as if they were fixed observations.\medskip

Existing data and code policies adopted by academic journals in the social sciences, like economics and management, offer only limited guidance for this situation. Most reproducibility standards require authors to document sources of randomness and provide code that controls them whenever possible \citep{vilhuber2019,koren2022,vilhuber2022}. These standards, however, largely predate the widespread adoption of generative AI and provide little guidance on how to report data generated by probabilistic black-box models. This lack of disclosure and guidance has important implications for the transparency and reproducibility of research, which are central concerns for academic journals \citep{fivsar2024reproducibility,loch2025reinforcing}.\medskip 

Our central thesis is simple: data generated by an LLM should be treated as draws from a distribution, not as fixed measurements. To develop this perspective, we make four contributions. First, we identify the main sources of variation across repeated model calls and explain the underlying mechanisms in a way that is accessible to the broad community of empirical researchers. Second, we distinguish sources of variation that researchers can control from those that remain under the control of model providers. Third, we propose an experiment that measures the consequences of these sources of variation in a typical task involving sentiment annotation of annual corporate filings. Fourth, we provide practical recommendations for (i) authors, to help them report the use of LLMs in their papers and prepare replication packages, and (ii) journal data editors, to help them assess these papers and their accompanying materials. We hope these contributions will stimulate discussion among researchers and journal (data) editors, and provide a foundation for the development of reporting and verification standards for the use of LLMs in empirical research.\medskip

% i think it's a good aim to help data editors, but we should keep in mind that there are about 10 of us, roughly, on the entire planet. Not sure we warrant a special mention here as a target group for the paper... maybe I would bracket the first (data) in data editor as well?

We complement the emerging literature on the reproducibility and reliability of LLM-generated data in empirical research. \citet{kim2026} shows that reliance on a single model run can produce unstable coefficients and misleading significance. The appropriate solution depends on whether model output is itself the research outcome or is an input for subsequent analysis. Related work by \citet{wang2025assessing} examines the consistency and reproducibility of LLM outputs across a range of finance and accounting tasks. They show that reproducibility varies across tasks, while simple aggregation over multiple model runs can substantially improve consistency. \citet{deKokMS2025} provides practical guidance for using LLMs in empirical research and is among the first papers to discuss the reproducibility implications of model-generated data. The paper recommends archiving prompts and raw model outputs, noting that commercial models may change or become unavailable over time.\medskip 

\section{How researchers use generative artificial intelligence}
\label{sec:researcher_use}

To document the rapid adoption of LLMs in economics and management research, we examine articles published between January 2020 and July 2026 in the 145 top-tier academic journals ranked 4 or 4* in the \href{https://charteredabs.org/academic-journal-guide}{Academic Journal Guide} (out of the 1,822 journals listed in the Guide). To identify articles using LLMs, we search the abstracts indexed by \href{https://www.openalex.org}{OpenAlex} for mentions of 212 LLMs (listed in Appendix \ref{sec:model_list}), as well as generic terms such as "LLM", "large language model", "Gen AI", "generative AI", and "generative artificial intelligence".\footnote{OpenAlex provides comprehensive metadata and abstracts for scholarly publications through an open bibliographic database and API. In the query, variations across the original terms are processed, too. For instance, plural forms are systematically tested and associated to the original term. Variations such as "GenAI" for "Gen AI" are also considered and searches are not case-sensitive (genAI counts just as well).}\medskip

Figure \ref{fig:term_freq} shows the share of articles in our sample of 105,464 publications that mention the above keywords. The increase is striking for all keyword groups. As of July 2026, nearly 6\% of papers published in leading economics and management journals mention LLMs in either their title or abstract. Although already substantial, this estimate is likely conservative because it is based solely on titles and abstracts rather than the full text of articles, and the 2026 value reflects publications indexed through July only. Moreover, published articles reflect research practices only with a lag of one to three years because of the time required for peer review and publication. Consistent with this view, a \href{https://www.linkedin.com/feed/update/urn:li:activity:7486740249776242688/}{July 2026 NBER survey of economists} found that 60\% reported LLMs to be an essential part of their research workflow.

\begin{figure}[!h]
    \centering
    \includegraphics[width=0.8\linewidth]{figures/term_adoption.pdf} \vspace{-5mm}
    \caption{\textbf{Use of LLMs in economics and business journals}. \small We plot the share of articles whose titles or abstracts feature a set of keywords associated with the use of LLMs. The sample consists of the 145 journals ranked 4 or 4* in the Academic Journal Guide ranking list. "Any model name" means that papers are flagged whenever they mention a specific LLM in their title or in their abstract. The total number of articles published in the 145 journals are 16,950 (2020), 17,153 (2021), 15,116 (2022), 14,602 (2023), 14,917 (2024), 16,414 (2025), and 10,312 (2026). \small}
    \label{fig:term_freq}
\end{figure}

Figure \ref{fig:open_closed} in the Appendix depicts the percentage of papers that mention specific LLMs, grouped by whether the models are closed-weight or open-weight.\footnote{Closed-weight models are proprietary and are typically accessed through the provider's infrastructure. Open-weight models can be downloaded and run locally by anyone with sufficiently powerful hardware. Open-source models are open-weight models for which the developers also release the training data, training code, and other components needed to reproduce the model.} This has important implications for researchers because open models often offer more control to users with respect to randomness of text generation. The plot shows that in recent years, researchers have been mainly using closed models (from OpenAI and Google mostly). If the trend continues and closed models do not allow for temperature tuning, then researchers will be working with highly random LLM outcomes in the future. Finally, Figure \ref{fig:provider_share} (also in the Appendix) reports the market share of LLM providers. While Google had the lead at the dawn of the decade with its BERT series of models, OpenAI now has the dominant market share among economics and business researchers.
\medskip

The increasing prevalence of LLMs documented in Figure \ref{fig:term_freq} raises a natural question: how are these models actually used in research in economics and management? We classify these uses into three main categories. The first is data generation and measurement. Researchers use models to classify text, annotate images, extract quantities, code qualitative evidence, and assign numerical scores. The outputs may become outcomes, explanatory variables, controls, or sample selection rules. This use creates the most direct reproducibility risk because LLM output becomes part of the data used in the analysis.\medskip

The second category involves simulated respondents and agents. Models can serve as substitutes for survey respondents or as agents in economic environments \citep{horton2023}. This approach can generate useful hypotheses, but it also raises questions about representativeness, embedded social biases, and the validity of replacing human behavior with model behavior \citep{bail2024}.\medskip

The third category is research assistance. Researchers use generative systems for programming, literature searches, writing, translation, and experimental design \citep{charness2023,korinek2023,korinek2025}. These uses create important disclosure and verification questions, but they usually leave artifacts that researchers can inspect. A human can review generated code, confirm a citation, or revise generated prose.\medskip

We focus on the first category because it has the most direct implications for research reproducibility. The central case is one in which model output becomes data. If a model determines a sentiment variable, identifies observations for inclusion, or measures an outcome, variation in its output can change the published result. Reproducibility therefore depends on the model and its serving environment as well as on the conventional data and code package.\medskip

We believe the issue is urgent because the adoption of LLMs in empirical research has been rapid. Many papers that rely on model-generated variables are already at advanced stages of the publication process, and many have already been published. Reporting norms are easier to establish before ad hoc reporting practices become entrenched.
% not sure i get the purpose of this last sentence.

\section{Where the randomness comes from}
\label{sec:randomness}

A language model produces text one token at a time. A token can be a word, part of a word, a number, or another unit used by the model. At each step, the model assigns a score to every candidate token. These scores are transformed into probabilities or used directly to select the next token.\medskip

Consider the incomplete sentence, ``\textit{John Maynard Keynes lived}'', and the particular LLM at hand forms predictions about the likelihood of the next token (word). Suppose the LLM assigns probability 0.70 to ``\textit{from}'', probability 0.15 to ``\textit{in}'', probability 0.10 to ``\textit{during}'', and probability 0.05 to all other candidates combined. LLMs typically apply a sampling procedure over those potential candidates, so as to make the model output less dull and repetitive. A greedy procedure would select ``\textit{from}'' because it has the highest score. The selected token is then added to the context, and the model calculates a new set of scores for the next token.\medskip

This sequential structure matters. A small difference at one step can alter every later step because later scores are conditional on the text already produced. Minor numerical variation can therefore lead to substantively different final answers.

\subsection{Deliberate sampling}

The most visible source of variation is deliberate sampling. Under sampling, the model draws the next token from the set of indexed candidate tokens $V$ (the vocabulary) rather than always choosing the token with the highest score. This design makes conversational output more varied and natural. Randomness is enforced as follows. The logits (scores) of all tokens are converted into probabilities for the next token $x_i$ through the following expression:
\begin{equation}
\label{eq:softmax}
P(x_i \mid \text{context}) \;=\; \frac{\exp(z_i/T)}{\sum_{j=1}^{|V|} \exp(z_j/T)},
\qquad T > 0,
\end{equation}
where $z_i$ is the logit score of token $i$, $|V|$ is the total number of possible tokens a model can choose from, and $T$ is the temperature. A high temperature flattens the distribution toward uniform (all words have the same probability to be chosen). A low temperature narrows it around the largest logit. In the limit, when the temperature is equal to zero (\textit{greedy decoding}), Equation \eqref{eq:softmax} is not defined and the model chooses the highest scoring token at each step. In the case of a tie, the token with the smallest token index $j$ is usually chosen.\footnote{In practice, temperature is rarely the only sampling parameter. It is typically combined with truncation rules that restrict the support of Equation \eqref{eq:softmax} before renormalizing: top-$k$ sampling retains only the $k$ highest-scoring tokens, while nucleus (top-$p$) sampling retains the smallest set of tokens whose cumulative probability exceeds $p$. Both discard the long tail of low-probability candidates, which is what keeps output coherent at higher temperatures; without them, the mass that a high $T$ shifts toward the uniform distribution is spread over a vocabulary of order $10^5$ tokens.}\medskip

Deliberate sampling is the source most directly controlled by the researcher when the model provider exposes generation parameters. Setting temperature to zero can substantially improve consistency and has no direct monetary cost.\footnote{However, there could be a cost in terms of model performance, see, e.g. \cite{du2025optimizing}.}\medskip

Unfortunately, as of mid-2026, several prominent reasoning models do not allow users to specify temperature. Providers such as OpenAI (with GPT-5.6) and Anthropic (with Sonnet 5, Opus 4.8 and Fable 5) have disabled the setting in favor of a reasoning intensity, which lets users tune the effort and cost related to their query while the sampling configuration is managed by the provider. In that case, the researcher cannot remove deliberate sampling directly. Figure \ref{fig:temperature} in the Appendix shows the timeline of model releases for four major LLM providers, along with their temperature policy. \medskip

A related control is the \textit{seed}. Several providers, including OpenAI, Google, Mistral, and Alibaba, expose a seed parameter to improve the reproducibility for some of their models. Temperature determines the sampling distribution in Equation \eqref{eq:softmax}, whereas the seed controls the pseudo-random draws from it. Seeds are irrelevant under greedy decoding and improve reproducibility only when the prompt, model, decoding settings, sampling implementation, and serving environment remain unchanged. Since hosted APIs do not guarantee these conditions, seeds should be reported when used, but they cannot ensure reproducibility.\medskip

\subsection{Silent model updates}

A provider may alter the model behind an unchanged public name. The endpoint, prompt, and user parameters can remain the same while the underlying weights or safety guardrails (e.g., Anthropic's Fable 5 between the beginning and the end of June 2026) can change. For instance, \citet{chen2024} document substantial variation in the behavior of GPT-4 over time. In one simple mathematical task involving the classification of numbers as prime or composite, reported accuracy fell from 84 percent to 51 percent within three months. Similarly, \citet{barrie2024} show that the performance of LLMs on repeated annotation tasks can vary substantially over time, making results difficult or even impossible to rerun and affecting downstream empirical findings. They argue that the precise source of this instability is unknown but may reflect either updates to the training data or undocumented modifications introduced by the model provider.\medskip

Software updates create familiar reproducibility problems, but model updates can be more severe. An older software version (e.g., Stata 15, released 2017) can often be archived and installed later. Though providers typically pin specific legacy models through their APIs, a proprietary model version may disappear completely. Even a dated model identifier may not reveal changes to system messages, filters, routing rules, or infrastructure.\medskip

Silent updates create both short term and long term risks. A model may change during data collection, so observations processed on different dates may not be comparable. It may also change after publication, preventing reproducibility verifiers from recovering the original outputs. Researchers using online services should therefore record the model identifier returned by every call, the date and time of execution, any provider fingerprint, and any available version information. These records do not ensure reproduction, but they make otherwise invisible changes easier to detect.

\subsection{Floating point rounding}

Computers store numbers using a finite number of bits. As a consequence, they cannot represent every real number: what is available is a grid of representable values, and any result falling between two grid points must be rounded to the nearest one. The crucial feature of this grid is that it is not evenly spaced. Near $1$, the grid points are about $2 \times 10^{-16}$ apart, so almost any quantity of practical interest can be represented to high accuracy. Near $10^{20}$, however, consecutive grid points are $16{,}384$ apart. The computer allocates its fixed budget of bits to the leading digits of a number, so the larger the number, the coarser the resolution with which it can be recorded.\medskip

This creates a problem whenever quantities of vastly different magnitude are combined. Adding $1$ to $10^{20}$ does not move the result far enough to reach the next grid point; the result falls between two representable values and is rounded back to where it started. The $1$ is silently discarded. Since rounding happens at every intermediate step, the order of operations matters:

\begin{verbatim}
> 1 + (1e20 - 1e20)
[1] 1
> (1 + 1e20) - 1e20
[1] 0
\end{verbatim}

In the first expression, the brackets signal to the programming language (R in this example) that $(1e20 - 1e20)$ should be computed before the sum should be taken. This implies that the two large terms cancel each other exactly, and the leading $1$ survives. In the second, the $1$ is absorbed into $10^{20}$ and lost before the subtraction ever takes place. The two expressions are mathematically identical, but one returns the correct answer and the other does not.%\footnote{Greedy decoding does not solve this: it is precisely the setting in which the arbitrariness is greatest, since sampling averages over the perturbed distribution while the argmax hands the entire outcome to whichever candidate the last digits happen to favour.}
\medskip

Actual model computations are vastly more complex, and numerical issues extend beyond addition. They include multiplication, reduced precision formats, accumulation rules, library implementations, and approximations used to accelerate matrix operations. The takeaway remains the same: finite precision arithmetic is not necessarily associative, and parallel execution can change the order of operations. Why might change the order?

\begin{enumerate}
\setlength\itemsep{-0.3em}
    \item One source is \textit{server load}. Online providers process many requests at once. The requests grouped into a batch can affect how operations are scheduled. The composition and size of the batch may differ across repeated calls. \medskip
    \item A second source is the \textit{hardware} on which the model runs. Different processors, drivers, numerical libraries, and precision settings can implement the same operation differently. Data centers may assign repeated requests to different machines or hardware configurations. They may also replace or upgrade this infrastructure without notifying users. \medskip
    \item A third source is \textit{model splitting}. Large models, because of their size, are often distributed across several processors. Partial calculations must then be combined. Differences in scheduling and communication can alter the order in which those calculations are performed.
\end{enumerate}

These mechanisms are not unique to language models. They are classical concerns in numerical and parallel computing. What distinguishes the predominant use of LLMs is that researchers often do not know and cannot control where the computations occur, which hardware performs them, or which other requests are processed at the same time.\medskip

Most changes in the final digits do not matter. If one candidate token has a clearly higher score than all others, small numerical differences will not alter the selected token. The problem arises when the leading candidates are nearly tied. In that case, even a tiny numerical change can reverse their order. Once a different token is selected, all later tokens are generated under a different context.\medskip

Recent empirical evidence confirms that this mechanism is consequential. \citet{ouyang2024} evaluate more than 800 code generation tasks and find that repeated requests frequently produce different outputs, even with a zero-temperature setting. \citet{he2025} identifies batch non-invariance as a key source of non-determinism, showing that server-side batching can alter numerical computations even under zero temperature. \citet{li2025} show that seemingly innocuous changes in the inference environment, including GPU type, GPU count, evaluation batch size, and numerical precision, can produce substantially different outputs by amplifying floating-point rounding errors. \citet{atil2024} systematically evaluate five LLMs across eight common tasks under nominally deterministic settings and find substantial variation across repeated runs, with accuracy differences reaching up to 15 percentage points.

\subsection{Expert routing}

Some LLMs use a mixture-of-experts architecture. Rather than using every part of the model for every token, the model sends each token to a small number of specialized components, known as experts. This allows the model to have greater overall capacity without a comparable increase in computing cost. However, each expert may be able to process only a limited number of tokens at a time. If too many tokens are sent to the same expert, some must either skip that expert or be sent elsewhere \citep{fedus2022switch}.\medskip

This creates a potential source of variation. Whether a token reaches its preferred expert can depend on which other users' requests happen to be processed in the same batch. As a result, two identical requests can, in principle, receive different outputs because they were processed alongside different requests. As part of a proof-of-concept adversarial attack, \citet{hayes2024} show that carefully chosen requests in the same batch can change a model's response to an otherwise unchanged input. The effect depends on how the model is implemented, but it is generally not visible to the user.\medskip

\subsection{What researchers can control}

Table~\ref{tab:sources} summarizes the main sources of variation. Local execution of open weights provides the strongest available control. Researchers can archive the weights, fix the batch size, preserve the software environment, record the hardware, and request deterministic numerical operations where supported. Exact reproduction may then be achievable, although it requires control of the complete serving stack rather than the weights alone. \medskip

\begin{table}[!h]
\centering
\small
\begin{tabular}{@{}L{2.9cm}L{7.4cm}C{1.8cm}C{1.1cm}C{1.6cm}@{}}
\toprule
\textbf{What differs between runs} & \textbf{In plain terms} & \textbf{Randomness survives $T{=}0$?} & \textbf{Control: local} & \textbf{Control: API} \\ 
\midrule 
\multicolumn{5}{@{}l}{\textit{\textcolor{sciblue}{Deliberate randomness}}}\\
\addlinespace[3pt]
Sampling & The model draws each word at random from its score distribution, instead of always taking the word with the highest score. & \textbf{no} & full  & for some models only \\
\addlinespace[6pt]
\multicolumn{5}{@{}l}{\textit{\textcolor{sciblue}{The model is not the same model}}}\\
\addlinespace[3pt]
Silent updates & The provider swaps the model behind an unchanged name. & yes & full & none \\  
\addlinespace[6pt]
\multicolumn{5}{@{}l}{\textit{\textcolor{sciblue}{Rounding: same arithmetic, different order, different results}}}\\
\addlinespace[3pt]
Server load & The prompt is computed alongside whatever other prompts arrived at the same moment, which changes how the arithmetic is divided up. %\emph{Main cause in practice.}
& yes & full & none \\
\addlinespace[2pt]
Hardware & A different chip, driver, or library version divides the same arithmetic differently. & yes & full & none \\
\addlinespace[2pt]
Model splitting & Very large models are cut across several chips, and the partial results are recombined in an order that varies. & yes & partial & none \\
\addlinespace[6pt]
\multicolumn{5}{@{}l}{\textit{\textcolor{sciblue}{A different part of the model ran}}}\\
\addlinespace[3pt]
Expert routing & Large models are a committee of experts, and each prompt is sent to a few of them. Places are limited, and other users' prompts can take them, so yours is handed to its second choice. & yes & partial & none \\
\bottomrule
\end{tabular}
\caption{\textbf{Differences between two runs of an identical prompt}. \small \textit{Local} means the model runs on the user's hardware thanks to open weights; \textit{API} means an online access point allowing the user to run the model on the provider's infrastructure.\label{tab:sources}}
\end{table}

Commercial services provide much less control, especially when temperature cannot be set. Researchers generally have no control over the hardware, batching behavior, or silent model updates. Open weight models, especially when deployed locally, circumvent this issue. They eliminate the dependence on the continued availability and policies of a private provider \citep{spirling2023}. Local execution nevertheless creates its own documentation requirements. The researcher must report the inference library, numerical precision, processor model, driver versions, batching settings, and any options related to deterministic execution.\medskip

\citet{yang2026} illustrates the importance of the serving stack. Across a large set of model and infrastructure combinations, locally served open models reproduced their own outputs much more consistently than proprietary systems. Even so, computers loading the same weights did not always agree. The same open model also displayed different consistency rates across cloud providers. In a financial portfolio application, average results were comparatively stable, but predicted directions changed for a meaningful share of firms. The central distinction is therefore not simply between open and proprietary models. It is between an execution environment that researchers can control and one they cannot observe or preserve.%\florian{it's not even clear that containerizing the app would solve this problem because it runs on different hardware (ie. floating point structures...)}

\section{Illustration with annual corporate filings}\label{sec:illustration}

We illustrate the practical consequences of model variation through a common measurement task. We follow the lines of \cite{tetlock2008more} and \cite{loughran2011liability}, who regress stock returns on proxies for sentiment. Here, sentiment is derived from LLM judgment on 10-K filings. The model classifies the filings' tone as positive, neutral, or negative. The experiment has two purposes. First, it measures output variation in a setting representative of applied work. Second, it separates deliberate sampling from residual variation where the available model settings permit this distinction.

\subsection{Experimental design}

The source documents are excerpts from Item 1 and Item 1A of annual filings submitted by S\&P 500 firms for 2025 on the SEC's EDGAR platform. Each excerpt is stored as a plain text file identified by ticker symbol. To limit API costs, the sample is restricted to the 100 firms with the shortest excerpts by token count. This restriction is purely for computational convenience, as the objective is to illustrate the existence of output variability rather than to estimate a population quantity. The following prompt was used on July 20, 2026: \medskip

{\ttfamily\raggedright Below is an excerpt from a company's 10-K filing (Item 1 Business and Item 1A Risk Factors). [EXCERPT] Classify the overall tone of this passage from a company's 10-K filing as positive, negative, or neutral for the company. Answer with exactly one word: positive, negative, or neutral.\par}

%For each firm, the prompt includes the filing excerpt and asks the model to classify the overall tone for the company as positive, negative, or neutral. The model must return \textit{precisely} one of those words. 

\medskip

The procedure is repeated 200 times with exactly the same prompt on the same model.%\medskip
This is carried out with three models: (1) OpenAI's GPT-5.6 Luna, (2) DeepSeek V4-Pro, and (3) Google's Gemma 4 26B A4B. While DeepSeek has not publicly disclosed a training-data or knowledge cutoff, GPT-5.6 was trained on data through February 2026, and Gemma 4 on data through January 2025. Accordingly, our sample may have been included in the training data available to GPT-5.6, but not in that available to Gemma 4. This difference could contribute to variation in the models’ outputs, although such analyses are beyond the scope of our paper.\medskip

For OpenAI, the reasoning level is set to "none". Temperature tuning is only available for models (2) and (3). For both models, we report results for both zero and nonzero temperatures. Sentiment scores are then digitized as follows: positive responses are coded as $1$, neutral responses as $0$, and negative responses as $-1$. For each repeated run, annual stock returns $r_{i,t}$ from the calendar year 2025 are regressed on the corresponding sentiment score $s_{i,t}$ (synchronous regression):

\begin{equation}
r_{i,t} = \alpha + \beta \times  s_{i,t} + e_{i,t}.
    \label{eq:sentiment}
\end{equation}

The regression produces a coefficient $\hat{\beta}$, a standard error and a $t$-statistic. It is the latter that we are interested in, because it captures both the sign and statistical significance of the coefficient.

\subsection{Results}

Figure \ref{fig:openai} shows the distribution of $t$-statistics associated with GPT-based sentiment scores. Because default temperature is nonzero, the generated scores are random, even if the set of outcomes is small (only three possible values for each firm in our experiment). The model reads the content of the filings and interprets them in different ways for each successive run.\footnote{The execution cost approximately \$49 for GPT-5.6 and completed in less than 100 minutes, with each firm taking between 40 seconds and 80 seconds to run (depending on filing length). Total token consumption was 274.2M for input and 0.1M in output (only one word output per call). For DeepSeek, the cost was \$3.7, split between 535.5M input tokens and 0.04M output tokens. Both models use prompt caching to avoid unnecessary token reading across the 200 runs.} \medskip

This first experiment shows that repeated classifications can produce very different regression statistics even though the filing excerpts, prompts, coding rules, and regression specification remain fixed. We notice that around 5\% of the $t$-statistics fall to the left of the 90\% confidence significance threshold. Hence, a researcher may in around 5\% of the cases report a statistically significant result, whereas in fact the related statistic passes the threshold only by luck. \medskip

\begin{figure}[!h]
\centering
\includegraphics[width=0.8\linewidth]{figures/gpt_56_luna_tstat_distribution.pdf} \vspace{-5mm}
\caption{\textbf{Distribution of regression statistics across repeated OpenAI model draws.} \small We show the distribution of test statistics across the 200 runs (Equation \eqref{eq:sentiment}). The two dashed vertical lines mark the thresholds for 90\% significance levels.}
\label{fig:openai}
\end{figure}

Figure \ref{fig:deepseek} repeats exactly the same exercise but in four different settings. First, we consider two alternative models: DeepSeek V4-Pro and Gemma-4-26B-A4B. The first is open-weight, but so large it can only be accessed from the provider's API. The second is also open-weight and its reasonable size makes it possible to run locally on a computer with 24 GB of RAM. Moreover, we carried out our tests with several temperature values: zero (to assess randomness under deterministic sampling) and one for DeepSeek (and 1.5 for Gemma).\footnote{These temperature values are mostly arbitrary and the value for Gemma-4 has been found by users on \href{https://www.reddit.com/r/LocalLLaMA/comments/1sg8r4l/gemma_4_seems_to_work_best_with_high_temperature/}{Reddit} to increase performance, compared to unit temperature.} \medskip

\begin{figure}[!h]
\centering
\begin{minipage}{0.495\linewidth}
\includegraphics[width=1\linewidth]{figures/deepseek_v4_pro_tstat_distribution.pdf} 
\end{minipage}
\begin{minipage}{0.495\linewidth}
\includegraphics[width=1\linewidth]{figures/Gemma_4_tstat_distribution.pdf} 
\end{minipage}
\vspace{-3mm}
\caption{\textbf{Distribution of regression statistics across repeated DeepSeek and Gemma-4 model runs.} \small We depict the distribution of test statistics from the regressions defined in Equation \eqref{eq:sentiment} over the 200 runs. The two dashed vertical lines mark the thresholds for the 90\% significance level. The bar color represents the chosen temperature. Light blue is used for $T>0$ (with frequencies reported in the left vertical axis) and green is used for $T=0$ (right vertical axis). } \label{fig:deepseek}
\end{figure}

The results exhibit markedly different degrees of dispersion. With unit temperature, the distribution of statistics based on DeepSeek-generated sentiment (left panel) is comparable to that obtained from OpenAI classifications in Figure \ref{fig:openai}. Under zero temperature, however, the distribution becomes much more concentrated but does not collapse to a single value. Instead, the estimates take only four distinct values (with one occurring only rarely). This residual variability arises because, among the 100 firms, two receive different sentiment classifications across repeated requests, yielding four possible combinations of outcomes. The result shows that setting $T=0$ removes deliberate sampling but, as explained in Section \ref{sec:randomness}, does not eliminate all sources of randomness in frontier models served through online APIs. By contrast, the distribution of statistics based on Gemma-4 deployed locally (right panel) is entirely concentrated at a single value under $T=0$. The residual variability disappears, indicating fully deterministic outputs under the conditions of our experiment.

\section{Practical guidance for reproducible science}

Economics has made substantial progress in research transparency during the past two decades. Data availability policies, registered reports, replication archives, journal reproducibility verification, and third party certification have improved the credibility of published evidence \citep{christensen2018,perignon2019,vilhuber2020,perignon2024a}. The AEA policy, the Social Science Data Editors template README, and the Data and Code Availability Standard (DCAS) provide widely used standards for replication packages \citep{vilhuber2019,vilhuber2022,koren2022}.\medskip

These systems require data, code, access instructions, provenance, and information about the computational environment. They do not yet provide sufficiently explicit guidance for model generated variables. The omission matters because conventional empirical pipelines already reproduce imperfectly at scale \citep{perignon2024b}. Generative models add another layer of complexity to an already demanding process.\medskip

% i think this is unnecessary:

% Subsection~\ref{sec:what-to-report} then classifies every item by where it belongs, in the paper or in the replication package, and by whether it is a must have or a good to have, and Subsection~\ref{sec:how-to-verify} describes the two checks a reproducibility verifier can run on the resulting package.

\subsection{What authors need to report}\label{sec:what-to-report}

There are two distinct places for reporting relevant details: (1) The paper contains what a reader needs in order to be able to judge the result: which model produced the data, under which settings, and how stable the output is. (2) The replication package carries what a reproducibility verifier needs to redo or bypass the model step: the exact machine readable artifacts, above all every model input and raw output before any post processing. Prompts belong here too, because prompts are code; a screenshot is not an adequate substitute, since it cannot be executed, searched, or compared.\medskip

Table~\ref{tab:reporting} classifies each relevant item along both dimensions and by priority. Items marked as \textit{must haves} are essentially free to provide, and no paper using model generated data should omit them. The \textit{good to have} items mostly involve repeated measurement or archiving; they entail a cost, which the experiment in Section~\ref{sec:illustration} makes concrete. %Appendix~\ref{app:reporting} discusses each group of items in detail and provides example reporting statements.
\medskip

\begin{table}[!h]
\centering
\small
\begin{tabular}{@{}L{5.2cm}L{4.4cm}L{4.6cm}@{}}
\toprule
\textbf{Item} & \textbf{In the paper} & \textbf{In the replication package} \\
\midrule
\multicolumn{3}{@{}l}{\textit{\textcolor{sciblue}{What the model was}}}\\
\addlinespace[3pt]
Model identification: name, exact version or checkpoint, provider, open or proprietary, access route, date when the model was queried & Must have & Must have (including hardware and serving stack for local execution) \\
\addlinespace[4pt]
Archived open weights or their cryptographic hash & --- & Good to have \\
\addlinespace[2pt]
Training data cutoff and its relation to the sample period & Must have (if available) & --- \\
\addlinespace[6pt]
\multicolumn{3}{@{}l}{\textit{\textcolor{sciblue}{How it was run}}}\\
\addlinespace[3pt]
Verbatim prompts, system messages, and examples & Must have (main prompt in an appendix) & Must have (prompts are code) \\
\addlinespace[4pt]
Generation parameters, distinguishing explicitly set values from provider defaults & Must have (headline settings, especially temperature) & Must have (complete configuration) \\
\addlinespace[6pt]
\multicolumn{3}{@{}l}{\textit{\textcolor{sciblue}{What it produced}}}\\
\addlinespace[3pt]
All model inputs and raw outputs from every run, before processing & --- & Must have\\
\addlinespace[4pt]
Variability report, number of draws, agreement measures, sensitivity of main estimates & Good to have (distribution of results for central claims) & --- \\
\addlinespace[6pt]
\multicolumn{3}{@{}l}{\textit{\textcolor{sciblue}{What it costed}}}\\
\addlinespace[3pt]
Cost statement: estimated token counts, monetary cost, runtime & --- & Must have \\
\bottomrule
\end{tabular}
\caption{\textbf{Reporting items for model generated variables}. \small Items are classified by where the information belongs and by priority.}\label{tab:reporting}
\end{table}

%The temperature setting is the only one that concerns practice rather than disclosure. Authors should use temperature zero for every output entering a reported result when the model permits it, and state explicitly when it does not. A seed should be fixed and reported where supported, but it supplements rather than replaces temperature control. Authors should nevertheless avoid claiming deterministic reproduction unless they have demonstrated it under an archived execution environment.\medskip

Cost is a legitimate concern, especially for large input data, but it does not justify ignoring variation. Repeated calls on a documented random subsample can provide an affordable assessment. Authors should explain the choice of repetition strategy and sample size in relation to the observed instability and the importance of the affected result.

\subsection{How data editors verify the results}\label{sec:how-to-verify}

The replication package column of Table~\ref{tab:reporting} supports two distinct checks that a data editor or reproducibility reviewer can run on a model based pipeline, each resting on a different group of items. The first check is \textit{reproduction from deposited outputs} using the archived raw outputs. The reproducibility reviewer skips model execution and uses those outputs to run the remaining parsing, cleaning, and statistical analysis. This check should reproduce the published results exactly, subject to conventional software and numerical considerations. This is why the raw output deposit is the most important item in the package: it works even after the model has changed, becomes too expensive to run, or disappears altogether.\medskip

The second is the \textit{replication of the results}, which rests on the deposited prompts, generation parameters, and execution instructions. Here the reviewer repeats the data generation itself. With a commercial service the new output may differ from the archive, so it should be read as a fresh model draw rather than an exact copy of the archived output. A variability report can supply the acceptance criterion: an estimate within the documented distribution supports the result, while one outside it signals a possible model update, infrastructure change, or underestimated variation. The README file should give an expected runtime and cost for regeneration. This is what verification teams and certification services need to plan and price their services. However, such analyses would be a new task for data editors as replications are typically not part of their mission.\medskip

\section{Conclusion}

The recent arrival of LLMs has opened many promising avenues for research. We focus on the role these models play as they become part of the data pipeline in empirical research. We point out that the practices adopted in the vast majority of current research imply severe issues for the reproducibility of results, and this paper is our attempt to start a serious conversation about remedying this situation.\medskip

% on purpose i don't say "this potentially dangerous situation". I mean "really dangerous".

We emphasize that there are various sources for the non-determinism of LLM output, even when prompted in the same way. We document deliberate sampling, silent updates, server batching, hardware differences, model splitting, and expert routing as sources. Some of those sources can be controlled by the user, while others remain out of the user's control like the deployment and order of computations in remote computing infrastructure with unknown, and continuously evolving, technical specifications. The use of open weight models has important advantages over proprietary models, not least because commercial model owners do not commit to preserving legacy models for future use. Nevertheless, even open weight models face challenges for reproducibility and require great care when describing the computational environment where the model was deployed. Taken together, this leads us to view LLM output as draws from a distribution rather than fixed measurements. \medskip

% i was wondering whether we should ask authors to report the full distribution of generated outputs? 

We propose a set of minimal requirements for researchers and journal editors to adopt. First, for researchers, we recommend setting temperature to zero whenever that option is available and the research objective does not require sampling. No other single decision removes as much variation at such low cost. Second, researchers should not describe temperature zero as deterministic. It removes deliberate sampling but leaves infrastructure and model variation in place. Third, where the model runs determines what the researcher can promise. A commercial service generally cannot support exact regeneration at any setting. Archived open weights on a controlled local stack provide the strongest route to exact reproduction, but only when the complete serving environment is documented and preserved. Containerization of local setups (via Docker, Singularity or similar) may be a promising solution for this purpose and should be explored. \medskip

Journals should insist on strict documentation requirements for all LLM usage, which should at a minimum encompass model name and version, any tuning parameter settings (like temperature), the full prompts in machine readable format, time and date the LLM was queried, and all generated output as raw data to be contained in the replication package. In general, all packages should envision a route to the full reproduction of the LLM outputs, starting from environment setup and prompts to input to the model (together with a monetary cost estimate), as well as a shortcut which takes the generated output as given.
\medskip

We share the profession's excitement about the promises held by LLMs for the future of research. It seems likely that many important findings will come from their judicious use and incorporation in research processes. Researchers typically state that \textit{standing on the shoulders of giants} is what enabled their particular advance. We want to caution that certain steps must be taken urgently to ensure that the foundations provided by this new class of giants are stable enough to build upon.

% LLM outputs should be treated as draws from a distribution rather than as fixed measurements. Their variation can arise from deliberate sampling, numerical rounding, server batching, hardware differences, model splitting, expert routing, and model updates. These sources differ in their magnitude, observability, and susceptibility to researcher control.\medskip

% Three lessons are central. First, set temperature to zero whenever that option is available and the research objective does not require sampling. No other single decision removes as much variation at such low cost. Second, do not describe temperature zero as deterministic. It removes deliberate sampling but leaves infrastructure and model variation in place. Third, where the model runs determines what the researcher can promise. A commercial service generally cannot support exact regeneration at any setting. Archived open weights on a controlled local stack provide the strongest route to exact reproduction, but only when the complete serving environment is documented and preserved.\medskip

% The practical response is straightforward. Authors should identify the model precisely, archive prompts and raw outputs, report all generation parameters, assess variation across repeated calls, and disclose the cost of regeneration. Replication packages should support both regeneration and bypass whenever possible.\medskip

% These practices also support emerging econometric methods for artificial intelligence generated variables. Valid correction and inference require information about measurement variation. A single unqualified model run does not provide that information. Repeated outputs and complete documentation do.\medskip

\bibliographystyle{apalike}
\renewcommand{\em}{\bfseries}
\bibliography{bib}
\nocite{ollion2024,palmer2024}
\clearpage
\appendix

\setcounter{figure}{0}
\renewcommand{\thefigure}{A\arabic{figure}}
\renewcommand{\theHfigure}{appendix.\arabic{figure}}
\begin{center} \Large
\textbf{APPENDIX}   
\end{center}

%\section{Supporting material}
\section{List of models}
\label{sec:model_list}

The 212 models considered in Section \ref{sec:researcher_use} are listed below, grouped by providers.

\begin{multicols}{2}
\begin{itemize} \small \setlength\itemsep{-0.4em}
  \item \textbf{01.AI}: Yi-34B, Yi-1.5, Yi-Large.
  \item \textbf{Academic}: BioBERT, FinBERT, ClinicalBERT, Legal-BERT.
  \item \textbf{AI21}: Jurassic-2, Jamba.
  \item \textbf{Alibaba}: Qwen, Qwen-7B, Qwen-VL, Qwen1.5, Qwen2, Qwen2.5, QwQ, Qwen2.5-Max, Qwen3, Qwen3.5, Tongyi Qianwen.
  \item \textbf{Allen Institute for AI}: SciBERT.
  \item \textbf{Amazon}: Amazon Titan, Amazon Nova.
  \item \textbf{Anthropic}: Claude 1, Claude 2, Claude Instant, Claude 3, Claude 3 Opus, Claude 3 Sonnet, Claude 3 Haiku, Claude 3.5 Sonnet, Claude 3.5 Haiku, Claude 3.7 Sonnet, Claude Opus 4, Claude Sonnet 4, Claude Opus 4.1, Claude Sonnet 4.5, Claude Haiku 4.5, Claude Opus 4.5, Claude Opus 4.6, Claude Opus 4.7, Claude Opus 4.8, Claude Sonnet 4.6, Claude Sonnet 5, Claude Fable 5.
  \item \textbf{Baichuan}: Baichuan, Baichuan2.
  \item \textbf{Baidu}: ERNIE 3.0, ERNIE Bot, ERNIE 4.0, ERNIE 4.5, ERNIE 5.0, ERNIE 5.1, Wenxin Yiyan.
  \item \textbf{BigCode}: StarCoder.
  \item \textbf{BigScience}: BLOOM language model.
  \item \textbf{ByteDance}: Doubao, Doubao 2.0, ByteDance Skylark.
  \item \textbf{Cohere}: Cohere Command, Cohere Command R, Cohere Command R+, Cohere Command A, Cohere Aya.
  \item \textbf{Databricks}: DBRX.
  \item \textbf{DeepSeek}: DeepSeek LLM, DeepSeek-Coder, DeepSeek-V2, DeepSeek-V3, DeepSeek-VL, DeepSeek-Math, DeepSeek-R1, DeepSeek-V3.1, DeepSeek-V3.2, DeepSeek-V4.
  \item \textbf{EleutherAI}: GPT-NeoX, GPT-J, Pythia.
  \item \textbf{Google}: Gopher, Chinchilla, LaMDA, Flan-T5, PaLM 2, Med-PaLM, Med-PaLM 2, Google Bard, Gemini 1.0, Gemini 1.5, Gemini 2.0, Gemini 2.5, Gemini 3, Gemini 3.5, Gemini 3.6, Gemma, Gemma 2, Gemma 3, Gemma 4, CodeGemma, BERT, ALBERT, ELECTRA, XLNet, mBERT.
  \item \textbf{Huawei}: Huawei PanGu.
  \item \textbf{Hugging Face}: DistilBERT.
  \item \textbf{iFlytek}: iFlytek Spark, SparkDesk.
  \item \textbf{Inspur}: Yuan 2.0.
  \item \textbf{Kunlun}: Skywork.
  \item \textbf{LMSYS/Academic}: Vicuna.
  \item \textbf{Meta}: OPT-175B, Galactica, LLaMA, Llama 2, Code Llama, Llama 3, Llama 3.1, Llama 3.2, Llama 3.3, Llama 4, Muse Spark, RoBERTa, XLM-R, XLM-RoBERTa.
  \item \textbf{Microsoft}: Phi-1, Phi-2, Phi-3, Phi-3.5, Phi-4, WizardLM, Orca, DeBERTa.
  \item \textbf{MiniMax}: MiniMax-Text, MiniMax-01, MiniMax M3.
  \item \textbf{Mistral}: Mistral 7B, Mixtral 8x7B, Mixtral 8x22B, Mistral Large, Mistral Large 2, Mistral Small, Mistral Medium, Mistral Nemo, Codestral, Mathstral, Pixtral, Ministral, Mistral Small 3, Mistral Medium 3, Magistral, Mistral Large 3, Mistral Small 4, Mistral Medium 3.5, Ministral 3.
  \item \textbf{Moonshot AI}: Kimi Moonshot, Kimi K1.5, Kimi K2, Kimi K2.5.
  \item \textbf{MosaicML}: MPT-7B.
  \item \textbf{NVIDIA}: Nemotron, Megatron-LM.
  \item \textbf{OpenAI}: GPT-2, GPT-3, InstructGPT, GPT-3.5, ChatGPT, GPT-4, GPT-4 Turbo, GPT-4V, GPT-4o, GPT-4o mini, GPT-4.1, GPT-4.5, OpenAI o1, o1-mini, OpenAI o3, o3-mini, o4-mini, GPT-5, GPT-5.1, GPT-5.2, GPT-5.4, GPT-5.5, GPT-5.6.
  \item \textbf{SenseTime}: SenseChat, SenseNova.
  \item \textbf{Shanghai AI Lab}: InternLM, InternLM2, InternLM2.5.
  \item \textbf{Stability AI}: StableLM.
  \item \textbf{Stanford/Academic}: Alpaca language model.
  \item \textbf{StepFun}: StepFun.
  \item \textbf{Tencent}: Hunyuan, Hunyuan-Large, Hunyuan-T1.
  \item \textbf{TII}: Falcon-40B, Falcon-180B.
  \item \textbf{xAI}: Grok-1, Grok-1.5, Grok 2, Grok 3, Grok 4, Grok 4.5.
  \item \textbf{Zhipu AI}: GLM-130B, ChatGLM, ChatGLM2, ChatGLM3, GLM-4, GLM-4.5, GLM-5, GLM-5.2, CogVLM.
\end{itemize}
\end{multicols}
\clearpage

\section{Open versus closed models}

\begin{figure}[!h]
    \centering
    \includegraphics[width=0.9\linewidth]{figures/open_vs_closed.pdf} \vspace{-5mm}
    \caption{\textbf{Open versus closed LLM models}. \small We chronologically represent the proportion of models that are mentioned in titles and abstracts, grouped according to whether they are open or closed. The total number of articles published in the 145 journals are 16,950 (2020), 17,153 (2021), 15,116 (2022), 14,602 (2023), 14,917 (2024), 16,414 (2025), and 10,312 (2026).}
    \label{fig:open_closed}
\end{figure}

\section{LLM provider market shares}

\begin{figure}[!h]
    \centering
    \includegraphics[width=0.9\linewidth]{figures/provider_share.pdf} \vspace{-5mm}
    \caption{\textbf{LLM Provider market shares}. \small We chronologically represent the share of models listed in titles and abstracts grouped by the organizations that develop them. In 2020, the lead from Google is due to the BERT model. Subsequently, academics have retrained BERT on domain-specific data, leading to specialized BERT versions, such as FinBERT, BioBERT, ClinicalBERT, Legal-BERT. These models fall under the umbrella of the "Academic" category in the legend.}
    \label{fig:provider_share}
\end{figure}

\section{Temperature timeline}

\begin{figure}[!h]
    \centering
    \includegraphics[width=1\linewidth]{figures/llm_temperature_timeline.pdf} \vspace{-7mm}
    \caption{\textbf{LLM temperature timeline (mid-2024 to mid-2026)}. \small We chronologically represent the release of models from Anthropic, DeepSeek, Google and OpenAI. Each model is associated with a policy regarding temperature: whether it can be user-specified (green), or whether it is simply ignored (orange), or prohibited (red). When temperature can be set, the range columns provide the range of possible values. The information was gathered from the providers' websites directly (for 75\% of points), and, when not available, from API calls, or from Wikipedia.}
    \label{fig:temperature}
\end{figure}


% \section{Example reporting statements}\label{app:reporting}

% This appendix expands on the reporting items summarized in Table~\ref{tab:reporting}. The four subsections follow the four groups of the table. For each item, we explain what authors should report and provide a template statement that can be adapted for the paper or the replication package.

% \subsection{What the model was}

% Authors should report the model name, exact version or checkpoint, provider, access route, and whether the model is proprietary or based on open weights. A suitable template for the README is:

% \begin{quote}
% \textit{Model generated variables were produced using [model name and exact identifier], accessed through [provider and interface] on [dates]. The model was [proprietary or open weights].}
% \end{quote}

% For local execution, authors should report the archived weights or their cryptographic hash, inference software, numerical precision, hardware, driver versions, batch size, and relevant serving settings:

% \begin{quote}
% \textit{For local execution, the archived weights have hash [hash], and inference used [software version], [numerical precision], [hardware], and batch size [value].}
% \end{quote}

% Authors should also state the model's reported training data cutoff and compare it with the sample period. If the sample predates the cutoff, the model may have encountered the documents or related information during training. This creates a contamination risk and may change the interpretation of the results. A suitable template for the paper is:

% \begin{quote}
% \textit{The provider reports a training data cutoff of [date]. The empirical sample covers [period]. Because [part or all] of the sample predates the cutoff, the model may have encountered the source documents or related information during training. The analysis should therefore be interpreted as [classification, extraction, prediction, or another task].}
% \end{quote}

% \subsection{How it was run}

% Prompts are code and should be included in the replication package as machine readable text files. Authors should preserve user prompts, system messages, examples, output schemas, and any instructions inserted automatically by a framework. Screenshots are not an adequate substitute because they cannot be executed, searched, or compared reliably. A suitable template for the README is:

% \begin{quote}
% \textit{All user prompts, system messages, examples, and output schemas are stored in folder [location].}
% \end{quote}

% In addition, authors should report every generation parameter they explicitly set and identify those left at provider defaults. Relevant settings include temperature, maximum output length, seed, reasoning level, batching, stopping rules, and structured output constraints. Defaults should not be treated as stable. A provider may change them across model versions. The README should therefore say that a parameter was left at its default rather than omit it:

% \begin{quote}
% \textit{Generation used temperature [value or unavailable], maximum output length [value], seed [value or unavailable], reasoning setting [value], batch size [value], and output constraint [description]. Parameters not listed as explicit values remained at provider defaults and may not be reproducible across versions.}
% \end{quote}

% \subsection{What it produced}

% Whenever legal, licensing, and confidentiality restrictions permit, authors should deposit every model input \textit{and} raw output. This deposit should occur before cleaning, parsing, or manual correction. It should include unsuccessful requests, invalid responses, and repeated draws not used in the final tables.\medskip

% Archived outputs allow reproducibility reviewers to bypass a model that has changed, become unavailable, or simply is too expensive to run. They also allow future researchers to apply alternative parsing and validation procedures.\medskip

% Any model generated variable that affects a substantive result should ideally be accompanied by a variability report, providing evidence about output stability. At minimum, authors should report the number of repeated calls, an agreement measure, the frequency of invalid outputs, and the sensitivity of the main estimates. The variability assessment can use the complete sample or a documented subsample. The selected subsample, randomization method, and seed should be preserved in the replication package. A suitable template for the README is:

% \begin{quote}
% \textit{We repeated the model generation step [number] times for [the complete sample or a random subsample of size]. Agreement across runs was [measure and value]. The main coefficient ranged from [value] to [value], changed sign in [share] of runs, and crossed the stated inference threshold in [share] of runs.}
% \end{quote}

% \subsection{What it costed}

% Authors should report the approximate number of input and output tokens, total monetary cost, runtime, rate limits, and local computational requirements. These quantities affect whether data editors and reproducibility reviewers can regenerate the data. A suitable template is:

% \begin{quote}
% \textit{The model generation stage processed approximately [number] input tokens and [number] output tokens. It cost [amount and currency] and required [runtime] as of [date]. Local execution, used [hardware] for [duration].}
% \end{quote}

\end{document}